\documentclass[10pt]{article}
\usepackage[margin=0.61in]{geometry}
\usepackage{amsmath,amssymb}
\usepackage[T1]{fontenc}
\usepackage[hidelinks]{hyperref}
\usepackage{microtype}

\newcommand{\Fock}{\mathcal F}
\newcommand{\Hs}{\mathcal H}

\title{The Singular Rotating-Wave Cutoff Limit\\
Converges in Norm Resolvent Sense}
\author{fa\_banach\_001 / agent\_lane\_00}
\date{17 August 2026}

\begin{document}
\maketitle
\vspace{-1.3em}
\small

\noindent\textbf{Status: literature already answered (affirmative under the
requested additional structure; no new theorem is claimed here).}

\section*{The 2022 question}

Lonigro constructs singular rotating-wave spin--boson Hamiltonians with
massive dispersion and form factor $f\in\Hs^r_{-s}\subset\Hs_{-s}$,
$1<s\le2$ and $s-1\le r\le1$.  After
Theorem~2, he notes that the genuinely renormalized cases are obtained from
regular cutoffs only in strong resolvent sense and asks whether this can be
improved, at least under additional assumptions \cite{Lon} (source PDF
page~10).

\section*{Direct later answer}

Hinrichs, Lampart, and Valent\'in Mart\'in treat
\[
 H_{\mathrm{SB},\Lambda}=A\otimes1+1\otimes d\Gamma(\omega)
   +B^*\otimes a(v_\Lambda)+B\otimes a^*(v_\Lambda).
\]
Their Example~2 identifies the rotating-wave approximation by
$B=\sigma_-$, hence $B^2=0$ \cite{HLV} (supporting PDF page~2).  Their
Theorem~2.5 (supporting PDF page~6) proves norm-resolvent convergence of the
renormalized regular approximants whenever the normal component is absent.
Consequently, for the pure rotating-wave interaction and
$(1+\omega)^{-1}v\in L^2$,
\[
 H_{\mathrm{SB},\Lambda}+B^*B E_\Lambda
    \xrightarrow[\Lambda\to\infty]{\ \mathrm{norm\ resolvent}\ }
 H_{\mathrm{SB}},
 \qquad E_\Lambda=\|\omega^{-1/2}v_\Lambda\|_2^2.
\]
Here $B^*B$ is the excited-state projection.  Thus this operator counterterm
is exactly a cutoff-dependent bare excitation-energy shift, matching the
source renormalization (with the coupling absorbed into $v$).

The source is massive.  On that domain, $f\in\Hs^r_{-s}\subset\Hs_{-s}$ for
$1<s\le2$ implies the later critical condition, and spectral cutoffs converge
in its weighted $b_2$ norm.  Hence the theorem covers the source's entire
singular range, including $s=2$.  The later authors explicitly state that
their result extends Lonigro's rotating-wave construction to arbitrary
coupling, massless fields, and norm-resolvent convergence.

\section*{Idea of the later proof}

For a 2-nilpotent interaction, the interior-boundary-condition construction
factorizes the Hamiltonian through
\[
 (1+G)\bigl(d\Gamma(\omega)+\lambda-T\bigr)(1+G^*),
 \qquad G^2=0,
\]
so $(1+G)^{-1}=1-G$.  The weighted $b_2$ topology makes $G$ and the
renormalized self-energy $T$ norm-continuous under cutoff removal.
Proposition~3.5 converts this into norm-resolvent continuity of the
factorized operator; Theorem~2.5 then adds the bounded spin term.  The
nilpotent case avoids the merely strong continuity of a nontrivial dressing
unitary, which is why the endpoint still has norm-resolvent convergence.

\section*{Conclusion and scope}

The open improvement question therefore has an affirmative answer for the
same rotating-wave model under its intrinsic 2-nilpotent structure, and in
fact with broader infrared and coupling scope.  This packet does not claim
that arbitrary nonnormal, nonnilpotent spin couplings enjoy the same limit.

\scriptsize
\begin{thebibliography}{9}
\bibitem{Lon} D.~Lonigro, \emph{Renormalization of spin--boson interactions
mediated by singular form factors}, arXiv:2210.15267 (2022).
\bibitem{HLV} B.~Hinrichs, J.~Lampart, and J.~Valent\'in Mart\'in,
\emph{Ultraviolet Renormalization of Spin Boson Models I. Normal and
2-Nilpotent Interactions}, arXiv:2502.04876 (2025).
\end{thebibliography}

\end{document}
